\section{Baccalauréat séries C/E — Session 2021}
\noindent\begin{minipage}{0.5\linewidth}\footnotesize Ministère des Enseignements Secondaires\\Office du Baccalauréat du Cameroun\end{minipage}%
\hfill\begin{minipage}{0.45\linewidth}\footnotesize\raggedleft Examen : \textbf{Baccalauréat} · Série : \textbf{C/E}\\
Épreuve : \textbf{Mathématiques} · Durée : \textbf{4\,h} · Coef. : \textbf{7\,(C) / 6\,(E)}\end{minipage}
\par\smallskip{\color{onbo}\rule{\linewidth}{1pt}}

\subsection*{Partie A — Évaluation des ressources \hfill (15 points)}

\noindent\textbf{Exercice 1.} \hfill\textit{(5,5 points pour la série C ; 4 points pour la série E)}\\[2pt]
\noindent\textbf{I -- \textit{(Série C exclusivement)}}\\
On considère la droite $(D)$ d'équation réduite $y=\dfrac{65}{16}x-\dfrac{5}{16}$ dans un repère orthonormé du plan.
\begin{enumerate}[label=\arabic*)]
\item Démontrer que $(D)$ passe par au moins un point $M$ dont les coordonnées sont des nombres entiers relatifs. \hfill\textit{(0,25 pt)}
\item Déterminer l'ensemble $E$ des points de $(D)$ à coordonnées entières. \hfill\textit{(0,75 pt)}
\item Déterminer les points de $(D)$ dont les ordonnées sont des entiers compris entre $-126$ et $134$. \hfill\textit{(0,5 pt)}
\end{enumerate}

\noindent\textbf{II -- } Soit un point $A(-2\,;\,1\,;\,1)$ et un vecteur $\vec n(1\,;\,-2\,;\,3)$ de l'espace $\mathcal E$ muni d'un repère orthonormé $(O\,;\,\vec\imath,\vec\jmath,\vec k)$.
\begin{enumerate}[label=\arabic*)]
\item Déterminer une équation du plan $(P)$ contenant le point $A$ et de vecteur normal $\vec n$. \hfill\textit{(0,5 pt)}
\item Donner une expression analytique de la réflexion de plan $(P)$. \hfill\textit{(1 pt)}
\end{enumerate}

\noindent\textbf{III -- } Le plan complexe est rapporté à un repère $(O\,;\,\vec u,\vec v)$. On considère la transformation $g$ du plan d'écriture complexe $Z'=\dfrac{1+i}{2}\,Z+1$.\\
$\Omega$ est le point d'affixe $1+i$, les points $A_n$ d'affixes $Z_n$.\\
$(Z_n)$ est la suite définie par : $Z_0=0$ et $Z_{n+1}=1+\dfrac{1+i}{2}\,Z_n$, pour tout entier naturel $n$.
\begin{enumerate}[label=\arabic*)]
\item Déterminer la nature et les éléments caractéristiques de $g$. \hfill\textit{(1 pt)}
\item Montrer que :
  \begin{enumerate}[label=\alph*)]
  \item Pour tout entier naturel $n$, les points $\Omega$, $A_n$ et $A_{n+4}$ sont alignés. \hfill\textit{(0,5 pt)}
  \item Pour tout entier naturel $n$, le triangle $\Omega A_n A_{n+1}$ est rectangle et isocèle. \hfill\textit{(1 pt)}
  \end{enumerate}
\end{enumerate}

\smallskip
\noindent\textbf{Exercice 2.} \hfill\textit{(4,5 points)}\\[2pt]
\noindent\textbf{I -- } Une urne contient $6$ boules indiscernables au toucher dont deux boules sont marquées $0$, trois boules sont marquées $\sqrt3$ et une boule est marquée $-\sqrt3$. On tire successivement et sans remise deux boules de cette urne.\\
On note $\lambda$ la variable aléatoire qui à chaque tirage associe la somme des nombres marqués sur les boules tirées.
\begin{enumerate}[label=\arabic*)]
\item Déterminer la loi de probabilité de $\lambda$. \hfill\textit{(0,75 pt)}
\item Calculer l'espérance mathématique et l'écart-type de $\lambda$. \hfill\textit{(0,75 pt)}
\end{enumerate}
\noindent\textbf{II -- } Le plan est muni d'un repère orthonormé direct $(O\,;\,\vec\imath,\vec\jmath)$.\\
$(\Sigma)$ est l'ensemble des points $M(X\,;\,Y)$ tels que $4X^2-Y^2=-4$.
\begin{enumerate}[label=\arabic*)]
\item Déterminer la nature et les éléments caractéristiques de $(\Sigma)$. \hfill\textit{(1 pt)}
\item[] $r$ est la rotation de centre $O$ et d'angle $-\dfrac{\pi}{6}$.
\item[2.] \begin{enumerate}[label=\alph*)]
  \item Donner l'expression analytique de $r$. \hfill\textit{(0,75 pt)}
  \item Déterminer une équation de l'ensemble $(\Sigma')$, image de $(\Sigma)$ par $r$. Déterminer la nature et les éléments caractéristiques de $(\Sigma')$. \hfill\textit{(0,5\,+\,1 pt)}
  \item Construire dans le repère $(O\,;\,\vec\imath,\vec\jmath)$ $(\Sigma)$ et $(\Sigma')$. \hfill\textit{(0,5 pt)}
  \end{enumerate}
\end{enumerate}

\smallskip
\noindent\textbf{Exercice 3.} \hfill\textit{(3,25 points pour la série C ; 4,75 points pour la série E)}\\[2pt]
On considère une fonction numérique $f$ définie sur $\R$ par $f(x)=\dfrac{x+2}{\mathrm e^{x}}$ et $(C)$ sa courbe représentative dans un repère orthonormé : unité sur les axes $2$\,cm.
\begin{enumerate}[label=\arabic*)]
\item \begin{enumerate}[label=\alph*)]
  \item Étudier les variations de $f$. \hfill\textit{(0,75 pt)}
  \item Déterminer une équation cartésienne de la tangente $(T)$ à $(C)$ au point d'abscisse $-1$. \hfill\textit{(0,25 pt)}
  \item Construire la courbe $(C)$ de $f$ et $(T)$ dans le même repère. \hfill\textit{(1 pt)}
  \end{enumerate}
\item \begin{enumerate}[label=\alph*)]
  \item Déterminer les constantes réelles $a$, $b$ et $c$ telles que la fonction $F$ définie sur $\R$ par $F(x)=\dfrac{ax+b}{\mathrm e^{x}}+cx$ soit une primitive de $f$. \hfill\textit{(0,75 pt)}
  \item Calculer $\displaystyle\int_{-1}^{0}f(x)\dd x$. \hfill\textit{(0,5 pt)}
  \end{enumerate}
\item \textbf{\textit{(E exclusivement)}}\\
On considère la fonction numérique $h$ définie sur $\R$ par $h(x)=f(-x)$ et $(C')$ sa courbe, et $(E)$ l'équation différentielle définie par : $y''-2y'+y=0$.
  \begin{enumerate}[label=\alph*)]
  \item Résoudre $(E)$. \hfill\textit{(0,75 pt)}
  \item Déterminer la solution de $(E)$ dont la courbe passe par le point $A(0\,;\,-1)$ et admet en ce point une tangente de coefficient directeur $1$. \hfill\textit{(0,75 pt)}
  \end{enumerate}
\end{enumerate}

\subsection*{Partie B — Évaluation des compétences \hfill (5 points)}

\noindent\textit{La figure ci-dessous représente le domaine d'un villageois nommé ABBA. Il a cultivé cette
année des carottes et des pastèques dans des portions comme l'indique la figure. Il a récolté le
même jour et a tout déversé dans un camion. Son fils KAM met en sac afin de vendre à raison de
\textbf{6\,800~F} le sac de pastèques et \textbf{3\,000~F} le sac de carottes, pour un total de \textbf{17} sacs.}\\[2pt]
\noindent\textit{À la fin de la vente, ABBA appelle KAM au téléphone pour savoir la recette obtenue. Avec des
problèmes de réseau, il le suit à peine et ne retient que : « la différence entre le prix de vente
total des carottes et des pastèques n'est que de \textbf{4\,000~F} ». Un sac de chaque type n'est pas vendu.}\\[2pt]
\noindent\textit{ABBA envisage de vendre une partie ou tout son vaste terrain à l'avenir. Dans cette zone, le m$^2$
coûte \textbf{2\,000~F}. Il confie ce projet à M.~KONG pour l'estimation de la valeur de ce terrain.
Celui-ci crée un repère indiqué sur la figure ci-dessous où l'unité sur l'axe des ordonnées est
\textbf{10~m} et \textbf{100~m} sur l'axe des abscisses. Les contours du terrain sont constitués de la droite $(AB)$,
la droite $(DB)$ et la ligne $(C)$. La droite $(L)$ représente la séparation de la portion exploitée
pour cultiver les pastèques de celle exploitée pour les carottes. KONG a réussi à trouver les
équations de $(C)$ et de $(L)$ qui sont respectivement $y=\dfrac{\mathrm e^{x}-\mathrm e^{-x}}{\mathrm e^{x}+\mathrm e^{-x}}$ et $y=\dfrac14 x$.}

\begin{illus}
\begin{tikzpicture}[scale=1.0,>=Latex]
  \def\xmax{4.3}
  % axes
  \draw[->] (-0.3,0) -- (\xmax,0) node[below right]{$x$};
  \draw[->] (0,-0.3) -- (0,1.6) node[left]{$y$};
  \foreach \x in {1,2,3,4} \draw (\x,0.04) -- (\x,-0.04) node[below,font=\scriptsize]{\x};
  \draw (0.04,0.5) -- (-0.04,0.5) node[left,font=\scriptsize]{0,5};
  \draw (0.04,1) -- (-0.04,1) node[left,font=\scriptsize]{1};
  % courbe (C) : tanh
  \draw[thick,onbink,domain=0:4,smooth,samples=80] plot (\x,{(exp(\x)-exp(-\x))/(exp(\x)+exp(-\x))});
  % droite (L) : y = x/4
  \draw[thick,onbblue,domain=0:4] plot (\x,{\x/4}) ;
  % points
  \fill (0,0) circle (1.3pt) node[below left,font=\scriptsize]{$A$};
  \fill (4,0) circle (1.3pt) node[below,font=\scriptsize]{$B$};
  \fill (4,1) circle (1.3pt) node[above right,font=\scriptsize]{$D$};
  % traits verticaux
  \draw[onbline] (4,0) -- (4,1);
  \node at (1.6,0.95) {\scriptsize pastèque};
  \node at (2.5,0.35) {\scriptsize carotte};
  \node[onbblue,font=\scriptsize] at (3.6,1.0) {$(L)$};
  \node[onbink,font=\scriptsize] at (1.0,0.85) {$(C)$};
\end{tikzpicture}
\fig{Domaine de ABBA : courbe $(C)$, droite $(L)$, contours $(AB)$, $(DB)$.}
\end{illus}

\begin{enumerate}[label=\textbf{Tâche \arabic*.},leftmargin=2.6em]
\item Combien coûtera ce terrain entier que ABBA souhaite vendre ? \hfill\textit{(1,5 pt)}
\item Combien aura ABBA s'il ne souhaite vendre que la portion réservée aux pastèques ? \hfill\textit{(1,5 pt)}
\item Aider ABBA à retrouver le nombre de sacs de chaque type des deux produits cultivés. \hfill\textit{(1,5 pt)}
\end{enumerate}
\par\smallskip{\footnotesize\itshape Présentation : 0,5 point.}

% ════════════════════════════════════════════════════
\corrige{de l'annale 2021}

\noindent\textbf{Partie A — Exercice 1.}

\smallskip
\noindent\textbf{I.\ 1)} Soit $(x\,;\,y)$ les coordonnées d'un point $M$ de $(D)$, avec $x,y\in\Z$. Alors
\[
y=\frac{65}{16}x-\frac{5}{16}\iff 16y=65x-5\iff 65x-16y=5 .
\]
Comme $65\wedge16=1$, le théorème de Bézout assure l'existence d'un couple $(x\,;\,y)\in\Z^2$ vérifiant $65x-16y=5$. Donc $(D)$ passe par au moins un point à coordonnées entières.

\noindent\textbf{2)} On résout $65x-16y=5$ dans $\Z^2$. Comme $65=16\times4+1$, on a $65-16\times4=1$, d'où $65\times5-16\times20=5$ : le couple $(5\,;\,20)$ est une solution particulière. Par soustraction,
\[
\begin{cases}65x-16y=5\\65\times5-16\times20=5\end{cases}\Rightarrow 65(x-5)-16(y-20)=0\Rightarrow 65(x-5)=16(y-20).
\]
Ainsi $16\mid 65(x-5)$ ; or $65\wedge16=1$, donc (Gauss) $16\mid x-5$. Il existe $k\in\Z$ tel que $x-5=16k$, soit $x=16k+5$. En reportant : $65\times16k=16(y-20)$, d'où $65k=y-20$, soit $y=65k+20$.

Réciproquement, $65(16k+5)-16(65k+20)=1040k+325-1040k-320=5$. Donc
\[
E=\{\,M(16k+5\,;\,65k+20),\ k\in\Z\,\}.
\]

\noindent\textbf{3)} On cherche $k\in\Z$ tel que $-126\le y\le134$ avec $y=65k+20$ :
\[
-126\le 65k+20\le134\iff -146\le65k\le114\iff -\tfrac{146}{65}\le k\le\tfrac{114}{65},
\]
soit $-2{,}25\le k\le1{,}75$. Les entiers conviennent : $k\in\{-2,-1,0,1\}$.
\[
\begin{array}{c|c|c}
k & x=16k+5 & y=65k+20\\\hline
-2 & -27 & -110\\
-1 & -11 & -45\\
0 & 5 & 20\\
1 & 21 & 85
\end{array}
\]
Les points cherchés sont $(-27\,;\,-110)$, $(-11\,;\,-45)$, $(5\,;\,20)$ et $(21\,;\,85)$.

\smallskip
\noindent\textbf{II.\ 1)} Tout plan d'équation cartésienne $ax+by+cz+d=0$ admet $\vec n(a\,;\,b\,;\,c)$ pour vecteur normal. Avec $\vec n(1\,;\,-2\,;\,3)$, $(P)$ a une équation de la forme $x-2y+3z+d=0$. Le point $A(-2\,;\,1\,;\,1)$ appartient à $(P)$ :
\[
-2-2\times1+3\times1+d=0\iff -1+d=0\iff d=1 .
\]
Donc $(P):\ x-2y+3z+1=0$.

\noindent\textbf{2)} Soit $M(x\,;\,y\,;\,z)$ d'image $M'(x'\,;\,y'\,;\,z')$ par la réflexion $S_{(P)}$. Alors $\overrightarrow{MM'}$ est colinéaire à $\vec n$ et le milieu $I$ de $[MM']$ appartient à $(P)$. Il existe $\lambda\in\R$ tel que $\overrightarrow{MM'}=\lambda\vec n$ :
\[
\begin{cases}x'-x=\lambda\\y'-y=-2\lambda\\z'-z=3\lambda\end{cases}\iff
\begin{cases}x'=x+\lambda\\y'=y-2\lambda\\z'=z+3\lambda.\end{cases}
\]
Le milieu $I\left(x+\dfrac{\lambda}{2}\,;\,y-\lambda\,;\,z+\dfrac{3\lambda}{2}\right)$ appartient à $(P)$ :
\[
\left(x+\tfrac{\lambda}{2}\right)-2(y-\lambda)+3\left(z+\tfrac{3\lambda}{2}\right)+1=0\iff 7\lambda+x-2y+3z+1=0,
\]
d'où $\lambda=-\dfrac17 x+\dfrac27 y-\dfrac37 z-\dfrac17$. En reportant dans le système, l'expression analytique de la réflexion de plan $(P)$ est :
\[
\begin{cases}
x'=\dfrac67 x+\dfrac27 y-\dfrac37 z-\dfrac17\\[4pt]
y'=\dfrac27 x+\dfrac37 y+\dfrac67 z+\dfrac27\\[4pt]
z'=-\dfrac37 x+\dfrac67 y-\dfrac27 z-\dfrac37 .
\end{cases}
\]

\smallskip
\noindent\textbf{III.\ 1)} $g$ a pour écriture $Z'=aZ+b$ avec $a=\dfrac{1+i}{2}\neq0$ et $b=1$ : c'est une similitude directe.
\begin{itemize}
\item Rapport : $k=|a|=\left|\dfrac{1+i}{2}\right|=\dfrac{\sqrt2}{2}$.
\item Angle : $\theta=\arg(a)=\arg(1+i)=\dfrac{\pi}{4}$.
\item Centre $\Omega_g$ d'affixe $\omega=\dfrac{b}{1-a}=\dfrac{1}{1-\frac{1+i}{2}}=\dfrac{2}{1-i}=\dfrac{2(1+i)}{2}=1+i$.
\end{itemize}
Donc $g$ est la similitude directe de centre $\Omega$ (affixe $1+i$), de rapport $\dfrac{\sqrt2}{2}$ et d'angle $\dfrac{\pi}{4}$.

\noindent\textbf{2.a)} On remarque que $Z_{n+1}=g(Z_n)$, donc $A_{n+1}=g(A_n)$. Pour l'angle orienté entre vecteurs successifs issus de $\Omega$ on a $\left(\overrightarrow{\Omega A_n}\,,\,\overrightarrow{\Omega A_{n+1}}\right)=\theta=\dfrac{\pi}{4}\ [2\pi]$. Par la relation de Chasles,
\[
\left(\overrightarrow{\Omega A_n}\,,\,\overrightarrow{\Omega A_{n+4}}\right)=\frac{\pi}{4}+\frac{\pi}{4}+\frac{\pi}{4}+\frac{\pi}{4}=\pi\ [2\pi].
\]
Cet angle valant $\pi$, les vecteurs $\overrightarrow{\Omega A_n}$ et $\overrightarrow{\Omega A_{n+4}}$ sont colinéaires : pour tout $n\in\N$, les points $\Omega$, $A_n$ et $A_{n+4}$ sont alignés.

\noindent\textbf{2.b)} Montrons que pour tout $n\in\N$, $\dfrac{z_{n+1}-z_n}{z_{n+1}-z_\Omega}=i$ (où $z_\Omega=1+i$). Comme $z_{n+1}=\dfrac{1+i}{2}z_n+1$ :
\begin{align*}
\frac{z_{n+1}-z_n}{z_{n+1}-z_\Omega}
&=\frac{\left(\frac{1+i}{2}z_n+1\right)-z_n}{\left(\frac{1+i}{2}z_n+1\right)-(1+i)}
=\frac{(1+i)z_n+2-2z_n}{(1+i)z_n+2-2-2i}\\
&=\frac{(-1+i)z_n+2}{(1+i)z_n-2i}
=\frac{\big[(-1+i)z_n+2\big]\times i}{\big[(1+i)z_n-2i\big]\times i}
=\frac{\big[(-1+i)z_n+2\big]\,i}{(-1+i)z_n+2}=i .
\end{align*}
Donc $\dfrac{z_n-z_\Omega}{z_{n+1}-z_\Omega}$ a pour module $1$ et pour argument $\dfrac{\pi}{2}$ : on en déduit
\[
\frac{\Omega A_n}{\Omega A_{n+1}}=1\quad\text{et}\quad \left(\overrightarrow{A_{n+1}A_n}\,,\,\overrightarrow{A_{n+1}\Omega}\right)=\frac{\pi}{2}.
\]
Ainsi $A_{n+1}\Omega=A_{n+1}A_n$ et l'angle en $A_{n+1}$ est droit : le triangle $\Omega A_n A_{n+1}$ est rectangle et isocèle en $A_{n+1}$.

\bigskip
\noindent\textbf{Exercice 2.}

\smallskip
\noindent\textbf{I.\ 1)} On tire successivement et sans remise deux boules parmi : deux marquées $0$, trois marquées $\sqrt3$, une marquée $-\sqrt3$. La somme $\lambda$ peut prendre les valeurs $-\sqrt3,\ 0,\ \sqrt3,\ 2\sqrt3$. Au premier tirage il y a $6$ boules, au second $5$. En détaillant chaque chemin de l'arbre :
\[
P(\lambda=0)=\frac26\times\frac15+\frac36\times\frac16+\frac16\times\frac36=\frac{2+3+3}{30}=\frac{8}{30}=\frac{4}{15},
\]
\[
P(\lambda=\sqrt3)=\frac26\times\frac36+\frac36\times\frac26=\frac{6+6}{30}=\frac{12}{30}=\frac{6}{15},
\]
\[
P(\lambda=-\sqrt3)=\frac26\times\frac16+\frac16\times\frac26=\frac{2+2}{30}=\frac{4}{30}=\frac{2}{15},
\]
\[
P(\lambda=2\sqrt3)=\frac36\times\frac26=\frac{6}{30}=\frac{3}{15}.
\]
\begin{center}
\begin{tabular}{c|cccc}
$\ell_i$ & $-\sqrt3$ & $0$ & $\sqrt3$ & $2\sqrt3$\\\hline
$P(\lambda=\ell_i)$ & $\dfrac{2}{15}$ & $\dfrac{4}{15}$ & $\dfrac{6}{15}$ & $\dfrac{3}{15}$\\
\end{tabular}
\end{center}
(On vérifie : $\frac{2+4+6+3}{15}=1$.)

\noindent\textbf{2)} Espérance :
\[
E(\lambda)=(-\sqrt3)\frac{2}{15}+0\cdot\frac{4}{15}+\sqrt3\cdot\frac{6}{15}+2\sqrt3\cdot\frac{3}{15}
=\frac{-2\sqrt3+6\sqrt3+6\sqrt3}{15}=\frac{10\sqrt3}{15}=\frac{2\sqrt3}{3}.
\]
Variance :
\[
V(\lambda)=\sum \ell_i^{\,2}P(\lambda=\ell_i)-\big(E(\lambda)\big)^2
=3\cdot\frac{2}{15}+0+3\cdot\frac{6}{15}+12\cdot\frac{3}{15}-\frac{12}{9}
=\frac{60}{15}-\frac{12}{9}=\frac{12}{3}-\frac{4}{3}=\frac{8}{3}.
\]
Écart-type :
\[
\sigma(\lambda)=\sqrt{V(\lambda)}=\sqrt{\frac{8}{3}}=\sqrt{\frac{24}{9}}=\frac{\sqrt{24}}{3}=\frac{2\sqrt6}{3}.
\]

\smallskip
\noindent\textbf{II.\ 1)} $4X^2-Y^2=-4\iff X^2-\dfrac{Y^2}{4}=-1$, soit $\dfrac{Y^2}{4}-X^2=1$. C'est une \textbf{hyperbole verticale} de la forme $\dfrac{Y^2}{a^2}-\dfrac{X^2}{b^2}=1$ avec $a^2=4$, $b^2=1$ ; on écrit ici $X^2-\dfrac{Y^2}{4}=-1$, d'où dans la forme $\frac{X^2}{1^2}-\frac{Y^2}{2^2}=-1$ : $a=1$, $b=2$, et $c^2=a^2+b^2=1+4=5$, $c=\sqrt5$.
\begin{itemize}
\item Centre : $O(0\,;\,0)$.
\item Asymptotes : $(a):\ Y=2X$ et $(a'):\ Y=-2X$.
\item Sommets : $S(0\,;\,2)$ et $S'(0\,;\,-2)$.
\item Foyers : $F(0\,;\,\sqrt5)$ et $F'(0\,;\,-\sqrt5)$.
\item Excentricité : $e=\dfrac{\sqrt5}{2}$.
\item Directrices : $(d):\ Y=\dfrac{4}{\sqrt5}=\dfrac{4\sqrt5}{5}$ et $(d'):\ Y=-\dfrac{4\sqrt5}{5}$.
\end{itemize}

\noindent\textbf{2.a)} L'image $M'$ d'affixe $z'=X'+iY'$ de $M$ d'affixe $z=X+iY$ par la rotation $r$ de centre $O$ et d'angle $-\dfrac{\pi}{6}$ vérifie $z'=\mathrm e^{-i\pi/6}z$ :
\[
X'+iY'=\left(\frac{\sqrt3}{2}-\frac12 i\right)(X+iY)
=\left(\frac{\sqrt3}{2}X+\frac12 Y\right)+i\left(-\frac12 X+\frac{\sqrt3}{2}Y\right).
\]
D'où l'expression analytique de $r$ :
\[
\begin{cases}X'=\dfrac{\sqrt3}{2}X+\dfrac12 Y\\[4pt]Y'=-\dfrac12 X+\dfrac{\sqrt3}{2}Y.\end{cases}
\]

\noindent\textbf{2.b)} On inverse : par addition convenable du système ci-dessus,
\[
X=\frac{\sqrt3}{2}X'-\frac12 Y',\qquad Y=\frac12 X'+\frac{\sqrt3}{2}Y'.
\]
Alors $M\in(\Sigma)\iff 4X^2-Y^2=-4$ :
\begin{align*}
4\left(\frac{\sqrt3}{2}X'-\frac12 Y'\right)^2-\left(\frac12 X'+\frac{\sqrt3}{2}Y'\right)^2&=-4\\
4\left(\frac34 X'^2-\frac{\sqrt3}{2}X'Y'+\frac14 Y'^2\right)-\left(\frac14 X'^2+\frac{\sqrt3}{2}X'Y'+\frac34 Y'^2\right)&=-4\\
\big(12X'^2-8\sqrt3\,X'Y'+4Y'^2\big)-\big(X'^2+2\sqrt3\,X'Y'+3Y'^2\big)&=-16\\
11X'^2+Y'^2-10\sqrt3\,X'Y'+16&=0 .
\end{align*}
Donc $(\Sigma'):\ 11X^2+Y^2-10\sqrt3\,XY+16=0$.

\noindent$(\Sigma')$ est une hyperbole : c'est l'image de l'hyperbole $(\Sigma)$ par la rotation $r$. Ses éléments caractéristiques sont les images par $r$ de ceux de $(\Sigma)$.
\begin{itemize}
\item Centre : $O(0\,;\,0)$ (image de $O$).
\item Première asymptote : image de $(a):Y=2X$, d'équation $Y=(-8+5\sqrt3)X$.\\
Deuxième asymptote : image de $(a'):Y=-2X$, d'équation $Y=(8+5\sqrt3)X$.
\item Sommets : image de $S(0\,;\,2)$ : $\left(\dfrac{\sqrt3}{2}\cdot0+\dfrac12\cdot2\,;\,-\dfrac12\cdot0+\dfrac{\sqrt3}{2}\cdot2\right)=(1\,;\,\sqrt3)$ ; par symétrie centrale, l'autre sommet est $(-1\,;\,-\sqrt3)$.
\item Foyers : image de $F(0\,;\,\sqrt5)$ : $\left(\dfrac{\sqrt5}{2}\,;\,\dfrac{\sqrt{15}}{2}\right)$ ; l'autre foyer est $\left(-\dfrac{\sqrt5}{2}\,;\,-\dfrac{\sqrt{15}}{2}\right)$.
\item Excentricité : $e=\dfrac{\sqrt5}{2}$ (inchangée par rotation).
\item Directrices : images de $(d)$ et $(d')$, d'équations approchées $Y\approx-0{,}58X+2{,}07$ et $Y\approx-0{,}58X-2{,}07$ (valeurs arrondies au centième).
\end{itemize}

\noindent\textbf{2.c)} La construction se fait en traçant l'hyperbole verticale $(\Sigma)$ (sommets $(0\,;\,\pm2)$, asymptotes $Y=\pm2X$), puis son image $(\Sigma')$ obtenue par la rotation de centre $O$ et d'angle $-\dfrac{\pi}{6}$.

\bigskip
\noindent\textbf{Exercice 3.}

\smallskip
\noindent\textbf{1.a)} $f(x)=\dfrac{x+2}{\mathrm e^{x}}$ est dérivable sur $\R$ (quotient de fonctions dérivables, $\mathrm e^{x}\neq0$) :
\[
f'(x)=\frac{1\cdot\mathrm e^{x}-(x+2)\mathrm e^{x}}{(\mathrm e^{x})^2}=\frac{(1-x-2)\mathrm e^{x}}{(\mathrm e^{x})^2}=\frac{(-x-1)\mathrm e^{x}}{(\mathrm e^{x})^2}=\frac{-x-1}{\mathrm e^{x}}.
\]
Comme $\mathrm e^{x}>0$, le signe de $f'(x)$ est celui de $-x-1$ : $f'(x)>0\iff x<-1$. Donc $f$ est strictement croissante sur $]-\infty\,;\,-1[$ et strictement décroissante sur $]-1\,;\,+\infty[$ ; elle admet un maximum en $x=-1$.

\noindent Limites : $f(-1)=\dfrac{-1+2}{\mathrm e^{-1}}=\mathrm e$ ; $\displaystyle\lim_{x\to-\infty}f(x)\overset{X=-x}{=}\lim_{X\to+\infty}(-X+2)\mathrm e^{X}=-\infty$ ; $\displaystyle\lim_{x\to+\infty}f(x)=0$ (croissances comparées).
\[
\begin{array}{c|ccccc}
x & -\infty & & -1 & & +\infty\\\hline
f'(x) & & + & 0 & - & \\\hline
& & & \mathrm e & &\\
f(x) & -\infty & \nearrow & & \searrow & 0
\end{array}
\]

\noindent\textbf{1.b)} La tangente $(T)$ en $x=-1$ a pour équation $y=f'(-1)(x+1)+f(-1)$. Or $f'(-1)=\dfrac{1-1}{\mathrm e^{-1}}=0$ et $f(-1)=\mathrm e$. Donc $(T):\ y=\mathrm e$.

\noindent\textbf{1.c)} La courbe $(C)$ part de $-\infty$, croît jusqu'au maximum $(-1\,;\,\mathrm e)$ où elle est tangente à la droite horizontale $(T):y=\mathrm e$, puis décroît vers l'asymptote horizontale $y=0$.

\noindent\textbf{2.a)} $F(x)=\dfrac{ax+b}{\mathrm e^{x}}+cx$ est primitive de $f$ ssi $F'=f$. Or
\[
F'(x)=\frac{a\,\mathrm e^{x}-(ax+b)\mathrm e^{x}}{(\mathrm e^{x})^2}+c=\frac{a-ax-b}{\mathrm e^{x}}+c=\frac{-ax+a-b+c\,\mathrm e^{x}}{\mathrm e^{x}}.
\]
On identifie avec $f(x)=\dfrac{x+2}{\mathrm e^{x}}$ :
\[
\begin{cases}-a=1\\a-b=2\\c=0\end{cases}\iff\begin{cases}a=-1\\b=-3\\c=0.\end{cases}
\]
Donc $F(x)=\dfrac{-x-3}{\mathrm e^{x}}$.

\noindent\textbf{2.b)}
\[
\int_{-1}^{0}f(x)\dd x=\Big[F(x)\Big]_{-1}^{0}=\left[\frac{-x-3}{\mathrm e^{x}}\right]_{-1}^{0}
=\frac{0-3}{\mathrm e^{0}}-\frac{1-3}{\mathrm e^{-1}}=-3-(-2\mathrm e)=-3+2\mathrm e.
\]

\noindent\textbf{3. (E exclusivement)}

\noindent\textbf{3.a)} L'équation caractéristique de $(E):\ y''-2y'+y=0$ est $r^2-2r+1=0$, soit $(r-1)^2=0$, d'où $r=1$ (racine double). La solution générale est
\[
y(x)=(\alpha x+\beta)\mathrm e^{x},\qquad (\alpha\,;\,\beta)\in\R^2.
\]

\noindent\textbf{3.b)} La solution $y_0$ vérifie $y_0(0)=-1$ (passage par $A(0\,;\,-1)$) et $y_0'(0)=1$ (coefficient directeur de la tangente). Or
\[
y_0'(x)=\alpha\,\mathrm e^{x}+(\alpha x+\beta)\mathrm e^{x}=(\alpha x+\alpha+\beta)\mathrm e^{x}.
\]
Donc
\[
\begin{cases}y_0(0)=\beta=-1\\y_0'(0)=\alpha+\beta=1\end{cases}\iff\begin{cases}\beta=-1\\\alpha=2.\end{cases}
\]
La solution cherchée est $y_0(x)=(2x-1)\mathrm e^{x}$.

\bigskip
\noindent\textbf{Partie B — Situation (ABBA).}

\smallskip
\noindent\textbf{Tâche 1.} On calcule d'abord l'aire $\mathcal A_1$ (en unités d'aire) de la surface comprise entre $(C)$, l'axe des abscisses et les droites $x=0$ et $x=4$ ($(C)$ est positive sur $[0\,;\,4]$). Comme $\dfrac{\mathrm e^{x}-\mathrm e^{-x}}{\mathrm e^{x}+\mathrm e^{-x}}=\dfrac{(\mathrm e^{x}+\mathrm e^{-x})'}{\mathrm e^{x}+\mathrm e^{-x}}$ :
\[
\mathcal A_1=\int_0^4\frac{\mathrm e^{x}-\mathrm e^{-x}}{\mathrm e^{x}+\mathrm e^{-x}}\dd x=\Big[\ln(\mathrm e^{x}+\mathrm e^{-x})\Big]_0^4
=\ln(\mathrm e^{4}+\mathrm e^{-4})-\ln 2=\ln\!\left(\frac{\mathrm e^{4}+\mathrm e^{-4}}{2}\right)\approx3{,}307\ \text{u.a.}
\]
L'unité valant $10$\,m (ordonnées) et $100$\,m (abscisses), une unité d'aire représente $10\times100=1000$\,m$^2$. Donc
\[
\mathcal A=1000\times3{,}307\approx3307\ \text{m}^2 .
\]
Le m$^2$ coûte $2\,000$\,F : le terrain entier coûte $3307\times2000=\mathbf{6\,614\,000}$\,F (valeur arrondie au millier).

\noindent\textbf{Tâche 2.} La portion « pastèque » est comprise entre $(C)$ (au-dessus) et $(L):y=\frac14 x$ (en dessous), sur $[0\,;\,4]$ :
\[
\mathcal A_2=\int_0^4\left(\frac{\mathrm e^{x}-\mathrm e^{-x}}{\mathrm e^{x}+\mathrm e^{-x}}-\frac14 x\right)\dd x
=\ln\!\left(\frac{\mathrm e^{4}+\mathrm e^{-4}}{2}\right)-\frac14\left[\frac{x^2}{2}\right]_0^4 .
\]
Or $\dfrac14\cdot\dfrac{4^2}{2}=2$, d'où
\[
\mathcal A_2=\ln\!\left(\frac{\mathrm e^{4}+\mathrm e^{-4}}{2}\right)-2\approx1{,}307\ \text{u.a.}
\]
soit $\mathcal A'=1000\times1{,}307\approx1307\ \text{m}^2$. Le montant de la portion pastèque est $1307\times2000=\mathbf{2\,614\,000}$\,F (valeur arrondie au millier).

\noindent\textbf{Tâche 3.} Soit $x$ le nombre de sacs de pastèques et $y$ celui de carottes. Au total $x+y=17$. Un sac de chaque type n'étant pas vendu, $(x-1)$ sacs de pastèques (à $6\,800$\,F) et $(y-1)$ sacs de carottes (à $3\,000$\,F) sont vendus, et la différence des recettes vaut $4\,000$\,F :
\[
6800(x-1)-3000(y-1)=4000 .
\]
On résout le système :
\[
\begin{aligned}
\begin{cases}x+y=17\\6800(x-1)-3000(y-1)=4000\end{cases}
&\iff\begin{cases}x+y=17\\34(x-1)-15(y-1)=20\end{cases}\\[4pt]
&\iff\begin{cases}x+y=17 & (1)\\34x-15y=39 & (2)\end{cases}
\end{aligned}
\]
$34\times(1)-(2)$ : $34y+15y=34\times17-39$, soit $49y=539$, d'où $y=11$. Puis $x=17-11=6$.

\noindent Donc ABBA a cultivé \textbf{6 sacs de pastèques} et \textbf{11 sacs de carottes}.
